\documentclass[12pt, a4paper]{article}
\usepackage[utf8]{inputenc}
\usepackage[T1]{fontenc}
\usepackage{amsmath, amssymb, amsthm, amsfonts}
\usepackage{geometry}
\usepackage{parskip} % Adds space between paragraphs
\usepackage{enumitem} % Better lists
\usepackage{fancyhdr} % For headers
\usepackage{bm} % Bold math for vectors
\usepackage{graphicx} % For graphics if needed

% Margins
\geometry{top=2.5cm, bottom=2.5cm, left=2.5cm, right=2.5cm}

% Theorem and Definition Styles
\theoremstyle{definition}
\newtheorem{definition}{Definition}
\newtheorem{example}{Example}
\newtheorem{theorem}{Theorem}
\newtheorem{lemma}{Lemma}
\newtheorem{remark}{Remark}
\newtheorem{corollary}{Corollary}

% Shortcuts for commonly used symbols
\newcommand{\R}{\mathbb{R}}
\newcommand{\E}{\mathbb{E}}
\newcommand{\Prob}{\mathbb{P}}
\newcommand{\Var}{\text{Var}}
\newcommand{\Cov}{\text{Cov}}
\newcommand{\iid}{\overset{\text{iid}}{\sim}}
\newcommand{\thetavec}{\underline{\theta}}
\newcommand{\xvec}{\underline{x}}
\newcommand{\Xvec}{\underline{X}}

\title{\textbf{Statistical Inference Class Notes}}
\author{Prateek N Jaiswal}
\date{November 27, 2025}

\begin{document}

\maketitle
\tableofcontents
\newpage

% -------------------------------------------------------------------
% PAGE 1 - 8: Introduction
% -------------------------------------------------------------------
\section{Introduction to Statistical Inference}

\subsection*{The Setup}
\begin{itemize}
    \item \textbf{Question:} What is the average height of a BSDS student?
    \item \textbf{Parameter of Interest:} The true average height (unknown constant).
    \item \textbf{Population:} All BSDS students.
    \item \textbf{Sample:} We observe the heights of the current batch of students.
\end{itemize}

\textbf{Goal:} To make inference about the parameter based on the sample at hand. To connect the gap between the sample and the population, we make use of statistical models.

\subsection*{Example 1: Student Heights}
\textbf{Assumption:}
The height of a BSDS student is $H$. We assume $H$ follows a Normal distribution:
$$ H \sim N(\mu, \sigma^2) $$

\textbf{Observations:}
We have observations $H_1, H_2, \dots, H_{62}$ which are independent and identically distributed (i.i.d.) copies of $H$.
$$ H_1, H_2, \dots, H_{62} \sim N(\mu, \sigma^2) $$

\textbf{Statistical Modeling:}
The average height of a BSDS student is the expectation:
$$ \E(H) = \mu $$
We want to estimate $\mu$ based on the sample $H_1, \dots, H_{62}$.

\subsection*{Example 2: LED Bulb Lifetime}
\textbf{Question:} What is the average lifetime of a Phillips LED bulb?
\begin{itemize}
    \item \textbf{Parameter of Interest:} Average lifetime.
    \item \textbf{Population:} All Phillips LED bulbs.
\end{itemize}

\textbf{Problem Scenario:}
If Phillips gives a one-year warranty on its bulbs, what proportion of bulbs need to be replaced?

\textbf{Experiment:}
Fix a time period $T$ (e.g., 1000 hours). Let 500 bulbs run for $T$ hours.
\begin{itemize}
    \item If a bulb is fused, we know exactly when it got fused (exact observation).
    \item If it is not fused by time $T$, we only know that it lasted for $> T$ hours.
\end{itemize}

\textbf{Modeling:}
Let the lifetime of a Phillips LED bulb be $L$.
$$ L \sim \text{Exponential}(\text{mean} = \lambda) $$
We observe $Y_1, \dots, Y_{500}$ where:
$$
Y_i = 
\begin{cases} 
L_i & \text{if } L_i \le T \quad (\text{Observed}) \\
T & \text{if } L_i > T \quad (\text{Truncated/Censored})
\end{cases}
$$
Here, $L_1, \dots, L_{500}$ are i.i.d. The observations $Y_i$ have the same distribution as $L$ truncated at $T$. The distribution of $Y$ involves $\lambda$ and $T$. From this, we can estimate $\lambda$.

\section{The Inference Framework}
In statistical inference:
\begin{enumerate}
    \item We have a \textbf{population} (often infinite).
    \item We have a \textbf{sample} from the population (finite).
    \item We postulate a certain \textbf{statistical model} on the population.
    \begin{itemize}
        \item Example: $H \sim N(\mu, \sigma^2)$
        \item Example: $L \sim \text{Exp}(\lambda)$
    \end{itemize}
    \item We also make certain assumptions on the sample collection mechanism (e.g., i.i.d., or truncated as in $\tilde{L}$).
\end{enumerate}

\subsection*{The Task}
Our task is to make inference about the \textbf{unknown parameters} of the population based on the sample.

\begin{definition}[\textbf{Parametric Model}]
When we assume the functional form of the population distribution to be known, except for the unknown parameters.
\end{definition}

\begin{definition}[\textbf{Nonparametric Model}]
We do not assume the form of the distribution to be known.
\begin{itemize}
    \item Example: $X \sim F$ where $f$ is continuous.
    \item We might only assume moments exist, e.g., $\E(X)$ exists, $\E(X^2)$ exists.
\end{itemize}
\end{definition}

\section{Parametric Models}
Throughout this course, we will focus on \textbf{parametric models}. We will assume we have observations:
$$ X_1, \dots, X_n \iid f_\theta $$
where $f_\theta$ is a known parametric form. The form of $f_\theta$ is known up to the parameter $\theta$. We know $f_\theta$ completely as soon as we know $\theta$.

\begin{example}[\textbf{Normal Distribution}]
$$ f_\theta = N(\mu, \sigma^2), \quad \theta = (\mu, \sigma^2) $$
The probability density function (pdf) is:
$$ f_{(\mu, \sigma^2)}(x) = \frac{1}{\sqrt{2\pi}\sigma} e^{-\frac{1}{2}\left(\frac{x-\mu}{\sigma}\right)^2}, \quad x \in \R $$
\end{example}

\begin{example}[\textbf{Exponential Distribution}]
Let $f_\theta = \text{Exponential}(\text{mean} = \lambda)$, where $\theta = \lambda$.
$$ f_\lambda(x) = \frac{1}{\lambda} e^{-\frac{x}{\lambda}}, \quad x > 0 $$
\end{example}

\begin{example}[\textbf{Binomial Distribution}]
Let $f_\theta = \text{Binomial}(n, p)$, where $\theta = p$.
$$ f_p(x) = \binom{n}{x} p^x (1-p)^{n-x}, \quad x = 0, 1, \dots, n $$
\end{example}

In general, we have a parametric model $X \sim f_\theta$, where $\theta$ is the unknown parameter.

% -------------------------------------------------------------------
% PAGE 10 - 22: Inference Goals, Parameters, Sufficiency
% -------------------------------------------------------------------
\section{Goals of Inference}
We have a parametric model $X_1, \dots, X_n \iid f_\theta$, where $\theta$ is the unknown parameter. We want to make inference about $\theta$.

\subsection{1. Estimation}
Estimate the value of $\theta$ based on $X_1, \dots, X_n$.
\begin{itemize}
    \item \textbf{Point Estimation:} We want a single value estimate.
    \item \textbf{Interval Estimation:} We want an interval where $\theta$ lies with high probability.
\end{itemize}

\subsection{2. Hypothesis Testing}
We want to check whether the $\theta$ belongs to some preconceived values that we have in mind.

\begin{example}
Let $X_1, \dots, X_n \sim \text{Binomial}(n, p)$.
\begin{itemize}
    \item \textbf{Estimation:} We may want to estimate $p$. A point estimator might be $\hat{p} = \frac{X}{n}$. An interval estimate might be $(\frac{X}{n} - c, \frac{X}{n} + c)$.
    \item \textbf{Testing:} We may want to test whether the value of $p$ equals 0.1.
    $$ H_0: p = 0.1 \quad \text{vs} \quad H_1: p \neq 0.1 $$
\end{itemize}
\end{example}

\section{The Parameter Space}
In our statistical model, the values of the parameters can vary within certain sets. These sets are known as the \textbf{Parameter Space}, denoted by $\Theta$.
\begin{itemize}
    \item \textbf{Normal:} $\Theta = \{(\mu, \sigma^2) : \mu \in \R, \sigma^2 > 0\}$
    \item \textbf{Exponential:} $\Theta = \{\lambda : \lambda > 0\}$
    \item \textbf{Binomial:} $\Theta = \{p : p \in [0, 1]\}$
\end{itemize}

\section{Data Reduction}
Given a sample $X_1, \dots, X_n \sim f_\theta$, the first step is often to summarize the data.
\begin{itemize}
    \item \textbf{Sample Mean:} $\bar{X} = \frac{1}{n}\sum_{i=1}^n X_i$.
    \item \textbf{Sample Variance:} $S^2 = \frac{1}{n-1}\sum_{i=1}^n (X_i - \bar{X})^2$.
    \item \textbf{Order Statistics:} $X_{(1)} = \min X_i$, $X_{(n)} = \max X_i$.
\end{itemize}

\begin{definition}[\textbf{Statistic}]
A statistic $T = T(\Xvec)$ is a measurable function of the sample. It is something that can be computed based on the sample alone.
\end{definition}

\section{Sufficiency}
\textbf{Idea:} Choose the summary statistic in such a way that \textbf{no information on $\theta$ is lost} in the process.

\begin{definition}[\textbf{Sufficient Statistic}]
We say that a statistic $T(\Xvec)$ is \textbf{sufficient} for $\theta$ if the conditional distribution of $\Xvec$ given $T(\Xvec)$ is \textbf{free of $\theta$}.
\end{definition}

\begin{example}[\textbf{Bernoulli Sufficiency}]
Let $X_1, \dots, X_n \iid \text{Bernoulli}(p)$. $T(\Xvec) = \sum X_i$.
$$ P_\theta(\Xvec = \xvec | T(\Xvec) = t) = \frac{p^t (1-p)^{n-t}}{\binom{n}{t} p^t (1-p)^{n-t}} = \frac{1}{\binom{n}{t}} $$
This is free of $p$, so $T$ is sufficient.
\end{example}

\begin{theorem}[\textbf{Factorization Theorem}]
Let $p_\theta(\xvec)$ be the joint pdf/pmf of the sample $\Xvec$. A statistic $T = T(\Xvec)$ is sufficient for $\theta$ if and only if we can factorize $p_\theta(\xvec)$ as:
$$ p_\theta(\xvec) = g_\theta(T(\xvec)) h(\xvec) $$
Where $g_\theta$ depends on $\theta$ and $T$, and $h(\xvec)$ is free of $\theta$.
\end{theorem}

% -------------------------------------------------------------------
% PAGE 33 - 60: Minimal Sufficiency
% -------------------------------------------------------------------
\section{Minimal Sufficiency}
A sufficient statistic reduces data. We want the one that gives the \textbf{most possible reduction}.

\begin{definition}[\textbf{Minimal Sufficient Statistic}]
A statistic $T(\Xvec)$ is minimal sufficient for $\theta$ if:
\begin{enumerate}
    \item $T(\Xvec)$ is sufficient.
    \item For any other sufficient statistic $S(\Xvec)$, $T(\Xvec)$ is a function of $S(\Xvec)$.
\end{enumerate}
\end{definition}

\begin{theorem}[\textbf{Lehmann-Scheff\'e Ratio Test}]
A statistic $T(\Xvec)$ is minimal sufficient for $\theta$ if and only if:
$$ \frac{p_\theta(\xvec)}{p_\theta(\mathbf{y})} \text{ is free of } \theta \iff T(\xvec) = T(\mathbf{y}) $$
\end{theorem}

\begin{example}[\textbf{Normal Mean}]
$X_i \iid N(\mu, 1)$. The likelihood ratio for $\xvec$ and $\mathbf{y}$ is free of $\mu$ iff $\sum x_i = \sum y_i$. Thus $T = \sum X_i$ is minimal sufficient.
\end{example}

\begin{example}[\textbf{Uniform}]
$X_i \iid \text{Uniform}(0, \theta)$. The ratio depends on indicators $\mathbb{I}(x_{(n)} < \theta)$. Independence from $\theta$ requires $x_{(n)} = y_{(n)}$. Thus $T = X_{(n)}$ is minimal sufficient.
\end{example}

% -------------------------------------------------------------------
% PAGE 61 - 84: Exponential Families & Order Statistics
% -------------------------------------------------------------------
\section{Exponential Families}
A family is a \textbf{one-parameter exponential family} if:
$$ f_\theta(x) = \exp(a(\theta)T(x) - b(\theta))c(x) $$
The statistic $\sum T(X_i)$ is minimal sufficient.

A family is a \textbf{$k$-parameter exponential family} if:
$$ f_\theta(x) = \exp\left( \sum_{j=1}^k a_j(\theta)T_j(x) - b(\theta) \right) c(x) $$
The vector $(\sum T_1(X_i), \dots, \sum T_k(X_i))$ is jointly minimal sufficient.

\section{Order Statistics}
The order statistics $X_{(1)} \le \dots \le X_{(n)}$ are always sufficient.
\subsection{Distributions}
For $X_i \iid f(x)$ with cdf $F(x)$:
\begin{itemize}
    \item \textbf{Minimum $X_{(1)}$:} $f_{(1)}(y) = n[1 - F(y)]^{n-1} f(y)$.
    \item \textbf{Maximum $X_{(n)}$:} $f_{(n)}(y) = n[F(y)]^{n-1} f(y)$.
    \item \textbf{General $X_{(j)}$:} $f_{(j)}(y) = \frac{n!}{(j-1)!(n-j)!} [F(y)]^{j-1} [1 - F(y)]^{n-j} f(y)$.
\end{itemize}

% -------------------------------------------------------------------
% PAGE 85 - 108: Ancillarity, Completeness, Basu
% -------------------------------------------------------------------
\section{Ancillary Statistics}
\begin{definition}
A statistic $S(\Xvec)$ is \textbf{ancillary} if its distribution is free of $\theta$.
\end{definition}
\begin{itemize}
    \item \textbf{Location Family:} $f(x-\theta)$. Range $R = X_{(n)} - X_{(1)}$ is ancillary.
    \item \textbf{Scale Family:} $\frac{1}{\theta} f(x/\theta)$. Ratios $X_{(i)}/X_{(j)}$ are ancillary.
\end{itemize}

\section{Completeness}
\begin{definition}[\textbf{Complete Statistic}]
A statistic $T$ is complete if $\E_\theta[g(T)] = 0$ for all $\theta$ implies $P(g(T)=0) = 1$ for all $\theta$.
\end{definition}

\begin{theorem}[\textbf{Basu's Theorem}]
If $T$ is a complete sufficient statistic and $S$ is an ancillary statistic, then $T$ and $S$ are \textbf{independent}.
\end{theorem}
\begin{example}
For $N(\mu, 1)$, $\bar{X}$ (complete sufficient) and $S^2$ (ancillary) are independent.
\end{example}

% -------------------------------------------------------------------
% PAGE 109 - 130: Estimation Methods
% -------------------------------------------------------------------
\section{Estimation Methods}
\subsection{Method of Moments (MOM)}
Equate population moments to sample moments: $\E[X^k] = \frac{1}{n}\sum X_i^k$. Solve for parameters.

\subsection{Maximum Likelihood Estimation (MLE)}
$$ \hat{\theta}_{MLE} = \arg\max_{\theta \in \Theta} L(\theta | \xvec) $$
where $L(\theta) = p_\theta(\xvec)$.

\section{Likelihood Principle}
If $L(\theta | \xvec) \propto L(\theta | \mathbf{y})$, inferences should be identical.

% -------------------------------------------------------------------
% PAGE 131 - 151: MLE Properties & Evaluation
% -------------------------------------------------------------------
\subsection{Invariance Property}
If $\hat{\theta}$ is the MLE of $\theta$, then $g(\hat{\theta})$ is the MLE of $g(\theta)$.

\subsection{MLE Examples}

\begin{example}[\textbf{Bernoulli MLE}]
Let $X_1, \dots, X_n \iid \text{Bernoulli}(p)$, $p \in [0, 1]$.
$$ L(p) = \prod_{i=1}^n p^{x_i} (1-p)^{1-x_i} = p^{\sum x_i} (1-p)^{n-\sum x_i} $$
Log-likelihood:
$$ \ell(p) = \left(\sum x_i\right) \ln p + \left(n - \sum x_i\right) \ln (1-p) $$
Differentiating with respect to $p$:
$$ \ell'(p) = \frac{\sum x_i}{p} - \frac{n - \sum x_i}{1-p} $$
Setting $\ell'(p) = 0$:
$$ (1-p) \sum x_i = p(n - \sum x_i) \implies \sum x_i - p \sum x_i = np - p \sum x_i $$
$$ \implies \hat{p}_{MLE} = \frac{\sum X_i}{n} = \bar{X} $$
\end{example}

\begin{example}[\textbf{Uniform MLE}]
Let $X_1, \dots, X_n \iid \text{Uniform}(0, \theta)$, $\theta > 0$.
$$ L(\theta) = \prod_{i=1}^n \frac{1}{\theta} \mathbb{I}(0 \le x_i \le \theta) = \frac{1}{\theta^n} \mathbb{I}(x_{(n)} \le \theta) \mathbb{I}(x_{(1)} \ge 0) $$
We want to maximize $L(\theta)$.
The function $\frac{1}{\theta^n}$ is strictly decreasing for $\theta > 0$.
However, the indicator $\mathbb{I}(x_{(n)} \le \theta)$ forces $\theta \ge x_{(n)}$.
To maximize a decreasing function subject to $\theta \ge x_{(n)}$, we choose the smallest possible valid value.
$$ \hat{\theta}_{MLE} = X_{(n)} $$
\end{example}

\subsection{Invariance Property of MLE}
Suppose we want to estimate a function of the parameter, say $\tau = g(\theta)$.
\begin{itemize}
    \item Example: $\theta = p$ (prob of success). We want to estimate $\tau = p^2$ (prob of two successes).
\end{itemize}

\begin{theorem}[\textbf{Invariance Property}]
If $\hat{\theta}$ is the MLE of $\theta$, then for any function $g(\cdot)$, the MLE of $g(\theta)$ is $g(\hat{\theta})$.
$$ \widehat{g(\theta)}_{MLE} = g(\hat{\theta}_{MLE}) $$
\end{theorem}

\subsection{Constrained Parameter Spaces}
The parameter space $\Theta$ plays an important role.
\begin{example}
Let $X \sim N(\mu, 1)$, but we know $\mu \ge 0$ (so $\Theta = [0, \infty)$).
The unconstrained maximum of the likelihood occurs at $\bar{X}$.
\begin{itemize}
    \item If $\bar{X} \ge 0$, then $\hat{\mu} = \bar{X}$ (inside $\Theta$).
    \item If $\bar{X} < 0$, the likelihood is increasing for $\mu < \bar{X}$ and decreasing for $\mu > \bar{X}$. Within $\Theta = [0, \infty)$, the maximum occurs at the boundary $\mu = 0$.
\end{itemize}
Thus, $\hat{\mu}_{MLE} = \max(0, \bar{X})$.
\end{example}

\textbf{Proof Sketch of Invariance:}
Define the induced likelihood for transformed parameter $\eta = g(\theta)$ as:
$$ L^*(\eta) = \sup_{\{\theta : g(\theta) = \eta\}} L(\theta) $$
We want to maximize $L^*(\eta)$.
$$ \sup_{\eta} L^*(\eta) = \sup_{\eta} \sup_{\{\theta : g(\theta)=\eta\}} L(\theta) = \sup_{\theta} L(\theta) = L(\hat{\theta}) $$
The maximum is achieved when $\eta = g(\hat{\theta})$.
Thus $\hat{\eta}_{MLE} = g(\hat{\theta})$.

\subsection{Numerical Methods for MLE}
In many cases (e.g., Gamma, Logistic regression), a closed-form solution for $\ell'(\theta) = 0$ does not exist. We use iterative optimization.

\textbf{1. Newton-Raphson Method}
To find the root of $\ell'(\theta) = 0$:
Start with an initial guess $\theta_0$.
Taylor expansion: $\ell'(\theta) \approx \ell'(\theta_0) + (\theta - \theta_0)\ell''(\theta_0)$.
Set to 0 to solve for next step:
$$ \theta_{n+1} = \theta_n - \frac{\ell'(\theta_n)}{\ell''(\theta_n)} $$
Here $-\ell''(\theta)$ is the \textbf{Observed Information}.

\textbf{2. Fisher Scoring}
Sometimes calculating or inverting the Hessian $\ell''(\theta)$ is unstable.
We can replace the observed second derivative with its expected value (Fisher Information).
$$ \E[\ell''(\theta)] = -I(\theta) $$
The update rule becomes:
$$ \theta_{n+1} = \theta_n + \frac{\ell'(\theta_n)}{I(\theta_n)} $$
(or in matrix form for multiparameter: $\theta_{n+1} = \theta_n + I^{-1}(\theta_n) \nabla \ell(\theta_n)$).

\section{Evaluating Estimators}

Let $T(\Xvec)$ be an estimator for $\theta$. The error is $T - \theta$.
We typically look at the \textbf{Mean Squared Error (MSE)}:
$$ MSE_\theta(T) = \E_\theta[(T - \theta)^2] $$

\textbf{Bias-Variance Decomposition:}
$$ MSE_\theta(T) = \Var_\theta(T) + [\text{Bias}_\theta(T)]^2 $$
where $\text{Bias}_\theta(T) = \E_\theta(T) - \theta$.

\subsection{Unbiasedness}
We would like to minimize MSE uniformly for all $\theta$. However, usually no single estimator minimizes MSE for \textit{all} $\theta$ (e.g., the estimator $T \equiv 17$ has zero MSE if $\theta=17$, but terrible MSE elsewhere).
To restrict the class of estimators, we often require Unbiasedness.

\begin{definition}[\textbf{Unbiased Estimator}]
An estimator $T$ is unbiased for $\theta$ if:
$$ \E_\theta(T) = \theta \quad \forall \theta \in \Theta $$
\end{definition}

\begin{definition}[\textbf{UMVUE}]
An estimator $T$ is a \textbf{Uniformly Minimum Variance Unbiased Estimator (UMVUE)} if:
\begin{enumerate}
    \item $\E_\theta(T) = \theta$ for all $\theta$.
    \item $\Var_\theta(T) \le \Var_\theta(S)$ for all $\theta$, for any other unbiased estimator $S$.
\end{enumerate}
\end{definition}

% -------------------------------------------------------------------
% PAGE 152 - 182: CRLB & UMVUE
% -------------------------------------------------------------------
\section{Cramer-Rao Lower Bound (CRLB)}
We want to find a lower bound on the variance of an unbiased estimator.

Let $X_1, \dots, X_n$ be a random sample with joint pdf/pmf $p_\theta(\xvec)$.
Let $T = T(\Xvec)$ be an estimator.
$$ \E_\theta[T(\Xvec)] = \int T(\xvec) p_\theta(\xvec) d\xvec $$
(Assuming the integral form; sum if discrete).

\textbf{Regularity Conditions:}
We assume the support of $p_\theta(\xvec)$ does not depend on $\theta$, and we can interchange differentiation and integration:
$$ \frac{\partial}{\partial \theta} \E_\theta[T] = \int T(\xvec) \frac{\partial}{\partial \theta} p_\theta(\xvec) d\xvec $$

Using the identity $\frac{\partial p_\theta}{\partial \ctheta} = p_\theta \frac{\partial \ln p_\theta}{\partial \theta}$:
$$ \frac{\partial}{\partial \theta} \E_\theta[T] = \int T(\xvec) \left( \frac{\partial \ln p_\theta(\xvec)}{\partial \theta} \right) p_\theta(\xvec) d\xvec = \E_\theta\left[ T(\Xvec) S_\theta(\Xvec) \right] $$
where $S_\theta(\Xvec) = \frac{\partial}{\partial \theta} \ln p_\theta(\Xvec)$ is the \textbf{Score Function}.

We also know that for the constant estimator 1:
$$ \int p_\theta(\xvec) d\xvec = 1 \implies \frac{\partial}{\partial \theta} \int p_\theta d\xvec = 0 \implies \E_\theta[S_\theta(\Xvec)] = 0 $$

\subsection{The Inequality}
Let $\tau(\theta) = \E_\theta[T]$. Then $\tau'(\theta) = \E[T S_\theta] = \text{Cov}(T, S_\theta)$ (since $\E[S_\theta]=0$).
By the Cauchy-Schwarz Inequality:
$$ [\text{Cov}(T, S_\theta)]^2 \le \Var_\theta(T) \Var_\theta(S_\theta) $$
$$ [\tau'(\theta)]^2 \le \Var_\theta(T) \E_\theta[(S_\theta)^2] $$
This gives the bound:
$$ \Var_\theta(T) \ge \frac{[\tau'(\theta)]^2}{\E_\theta[(\frac{\partial}{\partial \theta} \ln p_\theta(\Xvec))^2]} $$

\section{Fisher Information}

The denominator in the CRLB is called the Fisher Information.
For a single observation $X \sim f_\theta$:
$$ I(\theta) = \E_\theta\left[ \left( \frac{\partial}{\partial \theta} \ln f_\theta(X) \right)^2 \right] $$
Under regularity conditions (twice differentiable, interchange allowed), an alternative form is:
$$ I(\theta) = -\E_\theta\left[ \frac{\partial^2}{\partial \theta^2} \ln f_\theta(X) \right] $$

\subsection{Information in a Sample}
For an i.i.d. sample $X_1, \dots, X_n$:
$$ \ln p_\theta(\xvec) = \sum_{i=1}^n \ln f_\theta(x_i) $$
The information contained in the entire sample, denoted $I_n(\theta)$, is:
$$ I_n(\theta) = n I(\theta) $$
Thus, the Cramer-Rao Lower Bound for an unbiased estimator $W$ of $\theta$ (where $\tau(\theta)=\theta$) is:
$$ \Var_\theta(W) \ge \frac{1}{n I(\theta)} $$

\subsection{CRLB Examples}

\begin{example}[\textbf{Bernoulli}]
$X \sim \text{Bernoulli}(p)$.
$$ f(x) = p^x (1-p)^{1-x} \implies \ln f = x \ln p + (1-x) \ln (1-p) $$
$$ \frac{\partial \ln f}{\partial p} = \frac{x}{p} - \frac{1-x}{1-p} $$
$$ I(p) = \Var\left(\frac{X}{p} - \frac{1-X}{1-p}\right) = \frac{1}{p(1-p)} $$
CRLB for sample size $n$:
$$ \frac{1}{n I(p)} = \frac{p(1-p)}{n} $$
The variance of MLE $\hat{p} = \bar{X}$ is exactly $\frac{p(1-p)}{n}$. Thus, $\bar{X}$ attains the bound and is UMVUE.
\end{example}

\subsection{Multiparameter CRLB}
Let $\theta$ be a $k$-dimensional vector. Let $W$ be an unbiased estimator vector for $\theta$.
Let $\mathbf{I}(\theta)$ be the $k \times k$ Fisher Information Matrix:
$$ [\mathbf{I}(\theta)]_{ij} = -\E\left[ \frac{\partial^2}{\partial \theta_i \partial \theta_j} \ln p_\theta(\Xvec) \right] $$
The lower bound is given in terms of positive semi-definite (p.s.d.) matrices:
$$ \Var(W) \ge [\mathbf{I}(\theta)]^{-1} $$
Meaning $\Var(W) - [\mathbf{I}(\theta)]^{-1}$ is a p.s.d. matrix.

\subsection{Efficiency and Attainment}
An unbiased estimator $T$ that attains the CRLB is a UMVUE. However, the converse is not true; a UMVUE might exist but not attain the bound (e.g., if the bound is strictly unreachable).

\begin{example}[\textbf{Estimating Normal Variance}]
Let $X_i \iid N(\mu, 1)$. We want to estimate $\mu^2$.
$I(\mu) = 1$. The bound for estimating $g(\mu)=\mu^2$ involves $[g'(\mu)]^2$:
$$ \text{CRLB} = \frac{(2\mu)^2}{n \cdot 1} = \frac{4\mu^2}{n} $$
Consider the estimator $T = \bar{X}^2 - \frac{1}{n}$. It is unbiased for $\mu^2$.
$$ \Var(T) = \Var(\bar{X}^2) = \E[\bar{X}^4] - (\E[\bar{X}^2])^2 $$
Using $\bar{X} \sim N(\mu, 1/n)$, one can show:
$$ \Var(T) = \frac{4\mu^2}{n} + \frac{2}{n^2} $$
This is strictly greater than the CRLB $\frac{4\mu^2}{n}$. Yet, $T$ is the UMVUE (derived via completeness).
This shows the CRLB is not always sharp.
\end{example}

\begin{example}[\textbf{Uniform Case}]
$X_i \iid \text{Uniform}(0, \theta)$.
Support depends on $\theta$, so Regularity Conditions fail.
CRLB does not apply.
Variance of UMVUE ($X_{(n)}$ adjusted) is of order $O(1/n^2)$, whereas CRLB usually suggests $O(1/n)$.
\end{example}

\section{Characterization of UMVUE}

\begin{theorem}[\textbf{Covariance Method}]
Let $\mathcal{U}$ be the class of all unbiased estimators of $\theta$ (or $g(\theta)$).
Let $\mathcal{Z}$ be the class of all unbiased estimators of \textbf{zero} (i.e., $\E[Z] = 0$).
An estimator $T \in \mathcal{U}$ is UMVUE if and only if:
$$ \text{Cov}_\theta(T, Z) = 0 \quad \forall Z \in \mathcal{Z}, \forall \theta $$
\end{theorem}

\begin{proof}
\textbf{("If" part)}: Suppose $\text{Cov}(T, Z) = 0$ for all $Z$.
Let $S$ be any other unbiased estimator. Then $Z = T - S$ is an estimator of zero.
$$ \Var(S) = \Var(T - Z) = \Var(T) + \Var(Z) - 2\text{Cov}(T, Z) $$
$$ \Var(S) = \Var(T) + \Var(Z) \ge \Var(T) $$
Thus $T$ has minimum variance.

\textbf{("Only If" part)}: Suppose $T$ is UMVUE.
Consider $T_\lambda = T + \lambda Z$. This is unbiased.
$\Var(T_\lambda) = \Var(T) + \lambda^2 \Var(Z) + 2\lambda \text{Cov}(T, Z)$.
Minimize with respect to $\lambda$. If $\text{Cov}(T, Z) \neq 0$, we can choose $\lambda$ to make variance smaller than $\Var(T)$, contradicting UMVUE assumption.
Thus $\text{Cov}(T, Z)$ must be 0.
\end{proof}

\subsection{Properties of UMVUE}
\begin{enumerate}
    \item \textbf{Uniqueness:} The UMVUE is unique almost surely.
    \item \textbf{Linearity:} If $T_1$ is UMVUE for $g_1(\theta)$ and $T_2$ is UMVUE for $g_2(\theta)$, then $a T_1 + b T_2$ is UMVUE for $a g_1(\theta) + b g_2(\theta)$.
\end{enumerate}

\section{Rao-Blackwell Theorem}

This theorem gives a constructive way to improve estimators using sufficiency.

\begin{theorem}[\textbf{Rao-Blackwell}]
Let $T$ be an unbiased estimator of $\tau(\theta)$. Let $S$ be a sufficient statistic for $\theta$.
Define $T^* = \E[T | S]$. Then:
\begin{enumerate}
    \item $T^*$ is a valid statistic (free of $\theta$ due to sufficiency).
    \item $\E[T^*] = \tau(\theta)$ (Unbiased).
    \item $\Var(T^*) \le \Var(T)$ for all $\theta$.
\end{enumerate}
The inequality is strict unless $T$ is already a function of $S$.
\end{theorem}

\begin{proof}
Using the conditional variance formula:
$$ \Var(T) = \Var(\E[T|S]) + \E[\Var(T|S)] $$
$$ \Var(T) = \Var(T^*) + \E[\Var(T|S)] $$
Since $\Var(T|S) \ge 0$, we have $\Var(T) \ge \Var(T^*)$.
\end{proof}

\textbf{Note:} Conditioning on a sufficient statistic is crucial. Conditioning on a non-sufficient statistic might result in a variable that depends on $\theta$ (not a statistic) or doesn't improve variance optimally.

% -------------------------------------------------------------------
% PAGE 193 - 196: Lehmann-Scheffé
% -------------------------------------------------------------------
\section{Lehmann-Scheffé Theorem}

The Rao-Blackwell theorem gives us an improved estimator, but it doesn't guarantee uniqueness or optimality on its own unless we add the condition of completeness.

\begin{theorem}[\textbf{Lehmann-Scheffé}]
Let $T$ be an unbiased estimator of $\tau(\theta)$. Let $S$ be a \textbf{complete sufficient statistic} for $\theta$.
Then, $\phi(S) = \E[T | S]$ is the \textbf{unique UMVUE} of $\tau(\theta)$.
\end{theorem}

\textbf{Proof of Uniqueness:}
Suppose $T_1 = f_1(S)$ and $T_2 = f_2(S)$ are both unbiased estimators of $\tau(\theta)$ and are functions of the complete sufficient statistic $S$.
$$ \E_\theta[T_1 - T_2] = \tau(\theta) - \tau(\theta) = 0 \quad \forall \theta $$
Let $g(S) = T_1 - T_2$. Then $\E_\theta[g(S)] = 0$.
Since $S$ is complete, $P_\theta(g(S) = 0) = 1$.
Thus $T_1 = T_2$ almost surely.
Therefore, there is essentially only \textit{one} unbiased estimator that is a function of a complete sufficient statistic.

\subsection{Examples of UMVUE via Lehmann-Scheffé}

\begin{example}[\textbf{Bernoulli}]
$X_1, \dots, X_n \iid \text{Bernoulli}(p)$.
$S = \sum X_i$ is complete sufficient.
$\bar{X} = S/n$ is unbiased for $p$ and is a function of $S$.
Therefore, $\bar{X}$ is the UMVUE of $p$.
\end{example}

\begin{example}[\textbf{Poisson}]
$X_i \iid \text{Poisson}(\lambda)$. $\bar{X}$ is UMVUE for $\lambda$.
\end{example}

\begin{example}[\textbf{Uniform}]
$X_i \iid \text{Uniform}(0, \theta)$.
$X_{(n)}$ is complete sufficient.
$\E[X_{(n)}] = \frac{n}{n+1}\theta$.
Define $T = \frac{n+1}{n} X_{(n)}$.
$T$ is unbiased and is a function of the complete sufficient statistic.
Thus, $\frac{n+1}{n} X_{(n)}$ is the UMVUE of $\theta$.
\end{example}

\section{Tests of Hypotheses}

In hypothesis testing, we are interested to check whether $\theta \in \Theta_0$ or $\theta \in \Theta_1$, where $\Theta_0$ and $\Theta_1$ are non-overlapping subsets of the parameter space $\Theta$.

\begin{definition}[\textbf{Hypothesis}]
A hypothesis is a statement, assertion, or conjecture regarding the underlying statistical model (usually regarding the unknown parameter $\theta$).
\end{definition}

\begin{example}
$X_i \iid \text{Bernoulli}(p)$.
\begin{itemize}
    \item "The probability of success is 50\%" ($p=0.5$).
    \item "The probability is no more than 70\%" ($p \le 0.7$).
\end{itemize}
\end{example}

\begin{example}[\textbf{Non-parametric}]
$X_i \iid f$.
Hypothesis: "The distribution $f$ is Normal".
Hypothesis: "$f$ is symmetric about 0".
\end{example}

\subsection{Types of Hypotheses}

\textbf{1. Simple Hypothesis:}
If the hypothesis completely specifies the statistical model.
\begin{itemize}
    \item Example: $X_i \sim N(\mu, 1)$. $H: \mu = 0$. (Under $H$, the distribution is fully known as $N(0,1)$).
\end{itemize}

\textbf{2. Composite Hypothesis:}
If the hypothesis does not completely specify the distribution.
\begin{itemize}
    \item Example: $X_i \sim \text{Bernoulli}(p)$. $H: p \ge 0.7$. (Could be $p=0.7, p=0.8$, etc.)
    \item Example: $X_i \sim N(\mu, \sigma^2)$. $H: \mu = 0$. (Since $\sigma^2$ is unknown, the model is not fully specified).
\end{itemize}

\subsection{The Testing Framework}
We generally pit two hypotheses against each other:
\begin{itemize}
    \item $H_0$: The \textbf{Null Hypothesis} (usually denies change, represents status quo).
    \item $H_1$ (or $H_A$): The \textbf{Alternative Hypothesis}.
\end{itemize}

\begin{definition}[\textbf{Test Function}]
A test function is a rule that tells us whether to accept or reject a hypothesis based on the data $\Xvec$.
It is a map from the sample space $\mathcal{X}$ to the decision space \{Accept, Reject\}.
\end{definition}

\textbf{Partitioning the Sample Space:}
Usually, we partition $\mathcal{X}$ into two regions:
\begin{itemize}
    \item $\mathcal{R}$ (or $C$): The Rejection Region (or Critical Region). If $\xvec \in \mathcal{R}$, we reject $H_0$.
    \item $\mathcal{A}$: The Acceptance Region. If $\xvec \in \mathcal{A}$, we accept $H_0$.
\end{itemize}

\subsection{Randomized Tests}
Sometimes, especially with discrete data, we may need a "tie-breaker".
We toss a coin. If Heads, we reject; if Tails, we accept. This decision depends on an auxiliary random mechanism independent of data.

A general \textbf{Test Function} $\phi(\xvec)$ represents the probability of rejecting $H_0$ given data $\xvec$:
$$ 
\phi(\xvec) = 
\begin{cases} 
1 & \text{if } \xvec \in \mathcal{R} \quad (\text{Reject with prob 1}) \\
\gamma & \text{if } \xvec \in \mathcal{B} \quad (\text{Reject with prob } \gamma) \\
0 & \text{if } \xvec \in \mathcal{A} \quad (\text{Accept})
\end{cases}
$$
Here $\mathcal{B}$ is the randomization region.

\subsection{Decision Errors}

In a testing problem, there are four scenarios:

\begin{table}[h]
\centering
\begin{tabular}{|c|c|c|}
\hline
& \textbf{$H_0$ is True} & \textbf{$H_0$ is False ($H_1$ True)} \\
\hline
\textbf{Accept $H_0$} & Correct Decision & \textbf{Type II Error} \\
\hline
\textbf{Reject $H_0$} & \textbf{Type I Error} & Correct Decision \\
\hline
\end{tabular}
\end{table}

\textbf{Goal:} Ideally, we would like to make both errors as small as possible. However, simultaneously minimizing both is usually impossible (e.g., a test that always accepts $H_0$ has zero Type I error but maximal Type II error).

We focus on controlling the probabilities of these errors:
\begin{itemize}
    \item $P(\text{Type I Error}) = P(\text{Reject } H_0 | H_0 \text{ is True})$.
    \item $P(\text{Type II Error}) = P(\text{Accept } H_0 | H_0 \text{ is False})$.
\end{itemize}

\section{Power Function}

The quantity $\beta(\theta) = P_\theta(\text{Reject } H_0)$ is called the \textbf{Power Function} of the test.

\begin{itemize}
    \item If $\theta \in \Theta_0$ ($H_0$ is true), $\beta(\theta)$ is the probability of Type I Error.
    \item If $\theta \in \Theta_1$ ($H_1$ is true), $\beta(\theta)$ is the Power (probability of correctly rejecting).
    \item $P(\text{Type II Error}) = 1 - \beta(\theta)$ for $\theta \in \Theta_1$.
\end{itemize}

% -------------------------------------------------------------------
% PAGE 216 - 255: Optimal Tests
% -------------------------------------------------------------------
\section{Structure of Hypothesis Testing}

We typically have:
\begin{itemize}
    \item $H_0$: The Null Hypothesis.
    \item $H_1$ (or $H_A$): The Alternative Hypothesis.
\end{itemize}
These are pitted against each other. We accept one and reject the other based on the data.
Usually, $H_0$ denies change (status quo). We are reluctant to reject it unless we have strong evidence against it.
The alternative $H_1$ suggests a change. We accept it only if there is strong evidence in its support.

\subsection{Power Function Calculation}
Recall the randomized test function $\phi(\xvec)$:
$$ 
\phi(\xvec) = 
\begin{cases} 
1 & \text{if } \xvec \in \mathcal{R} \\
\gamma & \text{if } \xvec \in \mathcal{B} \\
0 & \text{if } \xvec \in \mathcal{A}
\end{cases}
$$
The Power Function $\beta(\theta)$ is the expected value of $\phi(\Xvec)$:
$$ \beta(\theta) = \E_\theta[\phi(\Xvec)] = 1 \cdot P_\theta(\Xvec \in \mathcal{R}) + \gamma \cdot P_\theta(\Xvec \in \mathcal{B}) + 0 \cdot P_\theta(\Xvec \in \mathcal{A}) $$
$$ \beta(\theta) = P_\theta(\mathcal{R}) + \gamma P_\theta(\mathcal{B}) $$
For a non-randomized test ($\gamma=0, \mathcal{B}=\emptyset$), this simplifies to just $P_\theta(\text{Reject } H_0)$.

\begin{example}[\textbf{Power in Uniform Distribution}]
Let $X_1, \dots, X_n \iid \text{Uniform}(0, \theta)$.
$H_0: \theta = 1$ vs. $H_1: \theta = 2$.
Suppose we use the test (based on $X_1$ only for simplicity):
$$ 
\phi(x) = 
\begin{cases} 
1 & \text{if } x_1 < 0.5 \\
\gamma & \text{if } 0.5 \le x_1 \le 1 \\
0 & \text{if } x_1 > 1
\end{cases}
$$
Power function $\beta(\theta) = \E_\theta[\phi(X)] = P_\theta(X_1 < 0.5) + \gamma P_\theta(0.5 \le X_1 \le 1)$.

\textbf{Case $\theta = 1$ ($H_0$):}
$$ P_{\theta=1}(X_1 < 0.5) = \int_0^{0.5} 1 dx = 0.5 $$
$$ P_{\theta=1}(0.5 \le X_1 \le 1) = \int_{0.5}^1 1 dx = 0.5 $$
$$ \beta(1) = 0.5 + 0.5\gamma $$

\textbf{Case $\theta = 2$ ($H_1$):} (pdf is $1/2$ on $(0,2)$)
$$ P_{\theta=2}(X_1 < 0.5) = \int_0^{0.5} \frac{1}{2} dx = 0.25 $$
$$ P_{\theta=2}(0.5 \le X_1 \le 1) = \int_{0.5}^1 \frac{1}{2} dx = 0.25 $$
$$ \beta(2) = 0.25 + 0.25\gamma $$
\end{example}

\section{Size and Level}

For a simple null hypothesis $H_0: \theta = \theta_0$, the probability of Type I error is exactly $\beta(\theta_0)$. This is called the \textbf{size} of the test.

For a composite null hypothesis $H_0: \theta \in \Theta_0$, the Type I error probability depends on $\theta$. We define the size as the worst-case error probability:
$$ \text{Size} = \sup_{\theta \in \Theta_0} \beta(\theta) $$

\begin{definition}[\textbf{Level $\alpha$ Test}]
A test is said to be of \textbf{level $\alpha$} if its size is at most $\alpha$.
$$ \sup_{\theta \in \Theta_0} \beta(\theta) \le \alpha $$
\end{definition}

\begin{definition}[\textbf{Unbiased Test}]
A test $\phi$ is unbiased if the power under the alternative is always greater than or equal to the size under the null.
$$ \inf_{\theta \in \Theta_1} \beta(\theta) \ge \sup_{\theta \in \Theta_0} \beta(\theta) $$
\end{definition}

\section{Optimal Tests}

We want to fix the size at $\alpha$ and maximize the power under the alternative.

\begin{definition}[\textbf{Most Powerful (MP) Test}]
A test $\phi$ is MP of level $\alpha$ for testing simple $H_0$ vs simple $H_1$ if:
\begin{enumerate}
    \item $\beta(\theta_0) \le \alpha$ (Size condition).
    \item $\beta(\theta_1) \ge \beta_{\phi'}(\theta_1)$ for any other level $\alpha$ test $\phi'$.
\end{enumerate}
\end{definition}

\begin{definition}[\textbf{Uniformly Most Powerful (UMP) Test}]
A test is UMP level $\alpha$ for $H_0$ vs composite $H_1$ if it is Most Powerful against \textit{every} $\theta \in \Theta_1$.
\end{definition}

\section{Neyman-Pearson Lemma}

This lemma gives the construction for the Most Powerful test for simple hypotheses.

\begin{theorem}[\textbf{Neyman-Pearson}]
Consider testing $H_0: \theta = \theta_0$ vs $H_1: \theta = \theta_1$ using observations with joint pdf/pmf $p_\theta(\xvec)$.
A test $\phi$ is Most Powerful of level $\alpha$ if and only if it has the form:
$$ 
\phi(\xvec) = 
\begin{cases} 
1 & \text{if } p_{\theta_1}(\xvec) > k p_{\theta_0}(\xvec) \\
\gamma & \text{if } p_{\theta_1}(\xvec) = k p_{\theta_0}(\xvec) \\
0 & \text{if } p_{\theta_1}(\xvec) < k p_{\theta_0}(\xvec)
\end{cases}
$$
where $k \ge 0$ and $\gamma \in [0,1]$ are constants determined by the size condition $\E_{\theta_0}[\phi(\Xvec)] = \alpha$.
\end{theorem}

\textbf{Remarks:}
\begin{itemize}
    \item \textbf{Sufficiency:} The MP test is sufficient; if an MP test exists, it must satisfy the Neyman-Pearson condition (except perhaps on a set of probability 0).
    \item \textbf{Uniqueness:} MP tests are not strictly unique, but they are unique if the boundary set $\{ \xvec : p_{\theta_1}(\xvec) = k p_{\theta_0}(\xvec) \}$ has probability 0 (which is true for continuous distributions). In that case, randomization ($\gamma$) is not needed.
    \item \textbf{Constants:} $k$ is usually positive. If $k=0$, we reject everywhere (trivial). If $k \to \infty$, we accept everywhere.
\end{itemize}

\begin{theorem}
An MP test is always unbiased.
\end{theorem}
\begin{proof}
(Sketch) We define sets $S^+ = \{p_1 > k p_0\}$ and $S^- = \{p_1 < k p_0\}$.
Using the inequality $(p_1 - k p_0)\phi \ge (p_1 - k p_0)\alpha$ (roughly), we can integrate to show:
$$ \int \phi (p_1 - k p_0) d\xvec \ge 0 $$
$$ \beta(\theta_1) - k \alpha \ge 0 \implies \beta(\theta_1) \ge k \alpha $$
(Detailed proof involves comparing against the constant test $\phi(\xvec) = \alpha$).
The result is $\beta(\theta_1) \ge \alpha$.
\end{proof}

\section{Monotone Likelihood Ratio (MLR)}

The Neyman-Pearson lemma applies to simple hypotheses. To find UMP tests for composite hypotheses (e.g., $H_1: \theta > \theta_0$), we often use the MLR property.

\begin{definition}[\textbf{Monotone Likelihood Ratio}]
A family of distributions $\{p_\theta(\xvec) : \theta \in \Theta\}$ has a Monotone Likelihood Ratio (MLR) in a statistic $T(\Xvec)$ if, for any $\theta_2 > \theta_1$, the ratio
$$ \frac{p_{\theta_2}(\xvec)}{p_{\theta_1}(\xvec)} $$
is a non-decreasing function of $T(\xvec)$.
\end{definition}

\subsection{UMP Tests via MLR}
\begin{theorem}
If the family has MLR in $T(\Xvec)$, then for testing $H_0: \theta \le \theta_0$ vs $H_1: \theta > \theta_0$, the UMP level $\alpha$ test is given by:
$$ 
\phi(\xvec) = 
\begin{cases} 
1 & \text{if } T(\xvec) > c \\
\gamma & \text{if } T(\xvec) = c \\
0 & \text{if } T(\xvec) < c
\end{cases}
$$
where $c$ and $\gamma$ are determined by $P_{\theta_0}(T(\Xvec) > c) + \gamma P_{\theta_0}(T(\Xvec) = c) = \alpha$.
\end{theorem}
Essentially, we reject $H_0$ for large values of the statistic $T$.

\begin{example}[\textbf{Exponential Family and MLR}]
For a one-parameter exponential family $f_\theta(x) = \exp(a(\theta)T(x) - b(\theta))c(x)$.
The likelihood ratio for $\theta_2 > \theta_1$ is:
$$ \frac{p_{\theta_2}}{p_{\theta_1}} = \exp\left( [a(\theta_2) - a(\theta_1)] T(\xvec) - n[b(\theta_2) - b(\theta_1)] \right) $$
If $a(\theta)$ is increasing in $\theta$, this ratio is increasing in $T(\xvec)$.
Thus, the family has MLR in $T(\Xvec)$.
\end{example}

\subsection{MLR in Exponential Families (Continued)}
Recall the one-parameter exponential family:
$$ f_\theta(x) = \exp(a(\theta)T(x) - b(\theta))c(x) $$
The minimal sufficient statistic is $T(\Xvec) = \sum T(x_i)$.

\begin{itemize}
    \item If $a(\theta)$ is strictly \textbf{increasing} in $\theta$, the family has MLR \textbf{increasing} in $T(\Xvec)$.
    The UMP test for $H_0: \theta \le \theta_0$ vs $H_1: \theta > \theta_0$ is:
    $$ \text{Reject if } T(\Xvec) > c $$
    \item If $a(\theta)$ is strictly \textbf{decreasing} in $\theta$, the family has MLR \textbf{decreasing} in $T(\Xvec)$.
    The UMP test for $H_0: \theta \le \theta_0$ vs $H_1: \theta > \theta_0$ is:
    $$ \text{Reject if } T(\Xvec) < c $$
\end{itemize}

\subsection{Properties of UMP Tests under MLR}
If the family has MLR in $T(\Xvec)$, then the power function $\beta_\phi(\theta)$ of the one-sided test (Reject for $T > c$) is \textbf{non-decreasing} in $\theta$.

Consequently:
\begin{enumerate}
    \item The test is of size $\alpha$ for the composite null $H_0: \theta \le \theta_0$.
    $$ \sup_{\theta \le \theta_0} \beta(\theta) = \beta(\theta_0) = \alpha $$
    \item The test is UMP for the composite hypotheses $H_0: \theta \le \theta_0$ vs $H_1: \theta > \theta_0$.
\end{enumerate}
Thus, under MLR, the "simple vs simple" Neyman-Pearson test automatically solves the "composite vs composite" one-sided problem.

\section{Test Statistics and Critical Values}

In many problems (even without MLR), we define a test using a Test Statistic $T(\Xvec)$ and a Cut-off (or Critical Value) $c$.
$$ \text{Reject } H_0 \iff T(\Xvec) > c \quad (\text{or } T(\Xvec) < c) $$
The critical value is determined by the null distribution of $T$ to satisfy the size condition:
$$ P_{\theta_0}(T(\Xvec) > c) = \alpha $$

\section{Two-Sided Hypotheses}
Consider the problem:
$$ H_0: \theta = \theta_0 \quad \text{vs} \quad H_1: \theta \neq \theta_0 $$
(i.e., $\theta > \theta_0$ or $\theta < \theta_0$).

A Uniformly Most Powerful (UMP) test generally does not exist for this problem.
\begin{itemize}
    \item The UMP test for $\theta > \theta_0$ rejects for large $T$.
    \item The UMP test for $\theta < \theta_0$ rejects for small $T$.
\end{itemize}
A test that rejects for large $T$ is usually "Uniformly \textbf{Least} Powerful" against the alternative $\theta < \theta_0$ (worse than a random guess). Since no single rejection region works for both sides, no UMP test exists.

\subsection{UMP Unbiased (UMPU) Tests}
To solve the two-sided problem, we restrict our search to Unbiased Tests.
Recall a test is unbiased if $\beta(\theta) \ge \alpha$ for all $\theta \neq \theta_0$.

\begin{theorem}
For a family with MLR in $T(\Xvec)$, the UMP Unbiased (UMPU) test for $H_0: \theta = \theta_0$ vs $H_1: \theta \neq \theta_0$ is of the form:
$$ 
\phi(\xvec) = 
\begin{cases} 
1 & \text{if } T(\xvec) < c_1 \text{ or } T(\xvec) > c_2 \\
\gamma_i & \text{if } T(\xvec) = c_i \quad (i=1,2) \\
0 & \text{if } c_1 < T(\xvec) < c_2
\end{cases}
$$
The constants are determined by two conditions:
\begin{enumerate}
    \item $\E_{\theta_0}[\phi(\Xvec)] = \alpha$ (Size condition)
    \item $\E_{\theta_0}[T(\Xvec) \phi(\Xvec)] = \alpha \E_{\theta_0}[T(\Xvec)]$ (Unbiasedness/Derivative condition)
\end{enumerate}
\end{theorem}

\begin{example}[\textbf{Normal Mean Two-Sided}]
$X_i \iid N(\mu, 1)$. Test $H_0: \mu = \mu_0$ vs $H_1: \mu \neq \mu_0$.
$T(\Xvec) = \bar{X}$.
The UMPU test rejects if $\bar{X} < c_1$ or $\bar{X} > c_2$.
Due to symmetry of the Normal distribution, the conditions simplify to symmetric critical values around $\mu_0$:
$$ c_1 = \mu_0 - z_{\alpha/2}/\sqrt{n}, \quad c_2 = \mu_0 + z_{\alpha/2}/\sqrt{n} $$
Reject if $|\bar{X} - \mu_0| > z_{\alpha/2}/\sqrt{n}$.
This matches the standard Z-test.
\end{example}

% -------------------------------------------------------------------
% PAGE 256 - 302: Large Sample Tests
% -------------------------------------------------------------------
\section{Likelihood Ratio Test (LRT)}

We now introduce general methods for constructing tests when UMP tests do not exist or are hard to find.

\subsection{The Statistic}
Consider testing $H_0: \theta \in \Theta_0$ vs $H_1: \theta \in \Theta \setminus \Theta_0$.
Let $\hat{\theta}_{MLE}$ be the unrestricted MLE over $\Theta$.
Let $\hat{\theta}_{0}$ be the restricted MLE over $\Theta_0$.

The Likelihood Ratio Statistic $\lambda(\xvec)$ is:
$$ \lambda(\xvec) = \frac{\sup_{\theta \in \Theta_0} L(\theta | \xvec)}{\sup_{\theta \in \Theta} L(\theta | \xvec)} = \frac{L(\hat{\theta}_0)}{L(\hat{\theta}_{MLE})} $$
Note that $0 \le \lambda(\xvec) \le 1$.
\begin{itemize}
    \item If $\lambda$ is close to 1, $H_0$ is supported.
    \item If $\lambda$ is close to 0, $H_0$ is rejected.
\end{itemize}

\subsection{Asymptotic Distribution (Wilks' Theorem)}
Finding the exact distribution of $\lambda(\xvec)$ can be difficult. We use the log-likelihood ratio.
Define the test statistic:
$$ -2 \ln \lambda(\xvec) = 2 [\ell(\hat{\theta}_{MLE}) - \ell(\hat{\theta}_0)] $$

\begin{theorem}[\textbf{Wilks}]
Under regularity conditions, as $n \to \infty$:
$$ -2 \ln \lambda(\Xvec) \xrightarrow{d} \chi^2_d $$
where the degrees of freedom $d$ is the difference in dimensions:
$$ d = \text{dim}(\Theta) - \text{dim}(\Theta_0) $$
(\# of free parameters in full model $-$ \# of free parameters under null).
\end{theorem}

\section{Consistency of Tests}

\begin{definition}[\textbf{Asymptotic Size}]
A sequence of tests $\phi_n$ has asymptotic size $\alpha$ if:
$$ \lim_{n \to \infty} \E_{\theta_0}[\phi_n(\Xvec)] = \alpha $$
\end{definition}

\begin{definition}[\textbf{Consistency}]
A test $\phi_n$ is consistent if for any fixed alternative $\theta \in \Theta_1$:
$$ \lim_{n \to \infty} \E_{\theta}[\phi_n(\Xvec)] = 1 $$
i.e., the Power goes to 1 as $n \to \infty$.
\end{definition}

The Likelihood Ratio Test (using the $\chi^2$ cutoff) is generally consistent.

\section{Wald's Test}

An alternative to the Likelihood Ratio Test that relies solely on the MLE and its asymptotic variance.

\subsection{The Setup}
Let $X_1, \dots, X_n$ be a sample from $p_\theta(\cdot)$.
We test $H_0: \theta = \theta_0$ vs $H_1: \theta \neq \theta_0$.
Let $\hat{\theta}_{MLE}$ be the unrestricted MLE.
Under regularity conditions, we know:
$$ \sqrt{n}(\hat{\theta}_{MLE} - \theta) \xrightarrow{d} N(0, I^{-1}(\theta)) $$
where $I(\theta)$ is the Fisher Information.

\subsection{The Statistic}
The Wald Test Statistic $T_W$ is based on the standardized distance between the MLE and the null value:
$$ T_W = \frac{(\hat{\theta}_{MLE} - \theta_0)^2}{\widehat{\Var}(\hat{\theta}_{MLE})} = n(\hat{\theta}_{MLE} - \theta_0)^2 I(\hat{\theta}_{MLE}) $$
Under $H_0$, as $n \to \infty$:
$$ T_W \xrightarrow{d} \chi^2_1 $$

\textbf{Decision Rule:}
Reject $H_0$ if $T_W > \chi^2_{1, \alpha}$ (the upper $\alpha$ critical value).

\subsection{Multiparameter Wald Test}
If $\theta$ is a $p$-dimensional vector, and we test $H_0: \theta = \theta_0$.
$$ \sqrt{n}(\hat{\theta}_{MLE} - \theta) \xrightarrow{d} N_p(\mathbf{0}, \mathbf{I}^{-1}(\theta)) $$
The test statistic becomes a quadratic form:
$$ T_W = n (\hat{\theta}_{MLE} - \theta_0)^T \mathbf{I}(\hat{\theta}_{MLE}) (\hat{\theta}_{MLE} - \theta_0) $$
Under $H_0$, $T_W \xrightarrow{d} \chi^2_p$.
We reject for large values of $T_W$.

\subsection{Alternative Forms and Nuances}
For the variance term in the denominator (or the Information matrix), we can use either:
\begin{enumerate}
    \item $I(\theta_0)$ (Expected information at null).
    \item $I(\hat{\theta}_{MLE})$ (Expected information at MLE).
    \item $J(\hat{\theta}_{MLE})$ (Observed information at MLE).
\end{enumerate}
Wald's test typically uses $I(\hat{\theta}_{MLE})$ or the Observed Information.
If we use $I(\theta_0)$, it is often called the Score Test (discussed later).
However, asymptotically they are equivalent.

\textbf{One-Sided Wald Test:}
For $H_0: \theta = \theta_0$ vs $H_1: \theta > \theta_0$.
Since the signed root $\sqrt{T_W}$ is asymptotically normal:
$$ Z_W = \frac{\sqrt{n}(\hat{\theta}_{MLE} - \theta_0)}{\sqrt{I^{-1}(\hat{\theta}_{MLE})}} \xrightarrow{d} N(0, 1) $$
We reject if $Z_W > z_\alpha$.
Conversely, for $H_1: \theta < \theta_0$, reject if $Z_W < -z_\alpha$.

\section{Testing Sub-vectors (Nuisance Parameters)}
Suppose $\theta = (\mu, \sigma^2)$. We want to test only $\mu = \mu_0$, treating $\sigma^2$ as a nuisance parameter.
Let $\theta = (\theta_1, \theta_2)$. Test $H_0: \theta_1 = \theta_{1,0}$.
We rely on the marginal asymptotic distribution of the MLE component $\hat{\theta}_1$.

If $\sqrt{n}(\hat{\theta} - \theta) \to N(\mathbf{0}, \Sigma)$, then the marginal distribution of the first component is:
$$ \sqrt{n}(\hat{\theta}_1 - \theta_1) \to N(\mathbf{0}, \Sigma_{11}) $$
where $\Sigma_{11}$ is the top-left block of the inverse information matrix $\mathbf{I}^{-1}(\theta)$.

The test statistic is:
$$ T_W = n (\hat{\theta}_1 - \theta_{1,0})^T [\Sigma_{11}(\hat{\theta})]^{-1} (\hat{\theta}_1 - \theta_{1,0}) $$
Under $H_0$, this follows $\chi^2_k$ where $k$ is the dimension of $\theta_1$.
This effectively "profiles out" the nuisance parameters by using the correct block of the inverse information matrix.

\section{Testing Functions of Parameters}
Often we care about a function $g(\theta)$ rather than $\theta$ itself.
Example: Test $H_0: g(\theta) = \gamma_0$ vs $H_1: g(\theta) \neq \gamma_0$.

\subsection{Delta Method}
If $\sqrt{n}(\hat{\theta} - \theta) \xrightarrow{d} N(0, \sigma^2)$, and $g$ is differentiable with $g'(\theta) \neq 0$:
$$ \sqrt{n}(g(\hat{\theta}) - g(\theta)) \xrightarrow{d} N(0, [g'(\theta)]^2 \sigma^2) $$

The Wald statistic for $g(\theta)$ is:
$$ T_W = \frac{n(g(\hat{\theta}) - \gamma_0)^2}{[g'(\hat{\theta})]^2 \widehat{\Var}(\hat{\theta})} $$
This converges to $\chi^2_1$.
Note: Wald's test is not invariant to transformations. Testing $\theta = 1$ vs testing $\ln \theta = 0$ can yield different results in finite samples.

\subsection{Multiparameter Delta Method}
Let $g: \R^p \to \R^q$ be a differentiable function.
Let $\nabla g(\theta)$ be the $q \times p$ Jacobian matrix of partial derivatives.
If $\sqrt{n}(\hat{\theta} - \theta) \xrightarrow{d} N_p(\mathbf{0}, \Sigma)$, then:
$$ \sqrt{n}(g(\hat{\theta}) - g(\theta)) \xrightarrow{d} N_q(\mathbf{0}, \nabla g(\theta) \Sigma \nabla g(\theta)^T) $$

\textbf{General Wald Statistic:}
$$ T_W = n (g(\hat{\theta}) - \gamma_0)^T \left[ \nabla g(\hat{\theta}) \widehat{\Sigma} \nabla g(\hat{\theta})^T \right]^{-1} (g(\hat{\theta}) - \gamma_0) $$
Under $H_0$, $T_W \xrightarrow{d} \chi^2_q$.

\begin{example}[\textbf{Linear Hypotheses}]
Consider testing $H_0: \mathbf{A}\theta = \mathbf{b}$.
Here $g(\theta) = \mathbf{A}\theta$, so $\nabla g(\theta) = \mathbf{A}$.
The variance term becomes $\mathbf{A} \mathbf{I}^{-1}(\theta) \mathbf{A}^T$.
$$ T_W = n (\mathbf{A}\hat{\theta} - \mathbf{b})^T (\mathbf{A} \widehat{\mathbf{I}^{-1}} \mathbf{A}^T)^{-1} (\mathbf{A}\hat{\theta} - \mathbf{b}) $$
Reject if $T_W > \chi^2_{q, \alpha}$.
\end{example}

\begin{example}[\textbf{Normal Example}]
$X_i \sim N(\mu, \sigma^2)$. $\theta = (\mu, \sigma^2)$.
Test $H_0: \mu = \mu_0$.
$g(\theta) = \mu$. $\nabla g = (1, 0)$.
Information Matrix for Normal is diagonal: $I(\mu, \sigma^2) = \text{diag}(1/\sigma^2, 1/2\sigma^4)$.
Inverse is $\text{diag}(\sigma^2, 2\sigma^4)$.
Variance of $\hat{\mu}$ is $\sigma^2/n$.
Wald statistic:
$$ T_W = \frac{(\bar{X} - \mu_0)^2}{\hat{\sigma}^2/n} = \frac{n(\bar{X} - \mu_0)^2}{\hat{\sigma}^2} $$
This is the squared Z-statistic (using MLE variance).
\end{example}

\textbf{Summary of Wald Strategy:}
Essentially, start with the limiting normal distribution of the entire vector $\hat{\theta}$. Find the marginal (or transformed) asymptotic distribution of the parameter of interest. Construct the quadratic form using the inverse of that specific variance block.

\section{Rao's Score Test}

The third major testing principle relies on the behavior of the gradient of the log-likelihood at the null hypothesis.

\subsection{The Score Function}
Let $X_1, \dots, X_n \iid f_\theta(x)$.
The log-likelihood is $\ell(\theta) = \sum_{i=1}^n \ln f_\theta(x_i)$.
The Score Function (or simply the Score) is defined as the gradient of the log-likelihood:
$$ S(\theta) = \ell'(\theta) = \sum_{i=1}^n \frac{\partial}{\partial \theta} \ln f_\theta(x_i) $$

\subsection{Properties of the Score}
Under regularity conditions:
\begin{enumerate}
    \item \textbf{Mean Zero:} $\E_\theta[S(\theta)] = 0$.
    $$ \E\left[\frac{\partial \ln f}{\partial \theta}\right] = \int \frac{1}{f} \frac{\partial f}{\partial \theta} f dx = \frac{\partial}{\partial \theta} \int f dx = 0 $$
    \item \textbf{Variance is Information:} $\Var_\theta(S(\theta)) = n I(\theta)$.
\end{enumerate}
By the Central Limit Theorem, since $S(\theta)$ is a sum of i.i.d. variables (individual scores $U_i = \frac{\partial}{\partial \theta} \ln f(X_i)$):
$$ \frac{1}{\sqrt{n}} S(\theta) \xrightarrow{d} N(0, I(\theta)) $$

\subsection{The Test Statistic}
Consider testing $H_0: \theta = \theta_0$ vs $H_1: \theta \neq \theta_0$.
If $H_0$ is true, the score evaluated at $\theta_0$, $S(\theta_0)$, should be close to its expected value of 0.
If $H_1$ is true (and the true $\theta$ is far from $\theta_0$), $S(\theta_0)$ will likely be large in magnitude (the likelihood slope is steep at $\theta_0$).

\textbf{Rao's Score Statistic:}
$$ T_S = \frac{[S(\theta_0)]^2}{\Var_{\theta_0}(S(\theta_0))} = \frac{[S(\theta_0)]^2}{n I(\theta_0)} $$
This statistic standardizes the score using the Fisher Information evaluated at the null.

\subsection{Asymptotic Distribution}
Under $H_0$, as $n \to \infty$:
$$ S(\theta_0) \approx N(0, n I(\theta_0)) $$
Therefore, the square of the standardized score follows a Chi-squared distribution:
$$ T_S \xrightarrow{d} \chi^2_1 $$

\textbf{Decision Rule:}
Reject $H_0$ if $T_S > \chi^2_{1, \alpha}$.
Note that this test statistic depends \textbf{only} on $\theta_0$ and the data. We do \textbf{not} need to compute the MLE $\hat{\theta}$. This is a major computational advantage.

\begin{example}[\textbf{Bernoulli Score Test}]
$X_i \iid \text{Bernoulli}(p)$. Test $H_0: p = p_0$.
Log-likelihood derivative:
$$ S(p) = \frac{\sum X_i}{p} - \frac{n - \sum X_i}{1-p} $$
Evaluated at $p_0$:
$$ S(p_0) = \frac{n \bar{X}}{p_0} - \frac{n(1-\bar{X})}{1-p_0} = \frac{n(\bar{X} - p_0)}{p_0(1-p_0)} $$
Fisher Information: $I(p_0) = \frac{1}{p_0(1-p_0)}$.
Score Statistic:
$$ T_S = \frac{\left[ \frac{n(\bar{X} - p_0)}{p_0(1-p_0)} \right]^2}{\frac{n}{p_0(1-p_0)}} = \frac{n (\bar{X} - p_0)^2}{p_0(1-p_0)} $$
This matches the standard Z-test for proportions ($\frac{(\hat{p}-p_0)^2}{p_0(1-p_0)/n}$).
\end{example}

\subsection{Score Test in Exponential Families}
For $f_\theta(x) = \exp(a(\theta)T(x) - b(\theta))c(x)$, the score is often linear in the sufficient statistic.
$$ S(\theta) = n a'(\theta) \bar{T} - n b'(\theta) $$
Thus, rejecting for large $|S(\theta_0)|$ is equivalent to rejecting for $\bar{T}$ far from its expected value under $H_0$.
Strategies for one-sided tests ($H_1: \theta > \theta_0$) rely on the sign of $S(\theta_0)$.
\begin{itemize}
    \item If $\E[S(\theta_0)] > 0$ for $\theta > \theta_0$, reject for large positive $S(\theta_0)$.
\end{itemize}

\subsection{Multiparameter Score Test}
Let $\theta$ be a $p$-dimensional vector.
The Score Vector is $\mathbf{S}(\theta) = \nabla \ell(\theta)$ ($p \times 1$ gradient).
Under $H_0: \theta = \theta_0$:
$$ \mathbf{S}(\theta_0) \xrightarrow{d} N_p(\mathbf{0}, \mathbf{I}_n(\theta_0)) $$
where $\mathbf{I}_n(\theta_0)$ is the Fisher Information Matrix.

\textbf{Statistic:}
$$ T_S = \mathbf{S}(\theta_0)^T [\mathbf{I}_n(\theta_0)]^{-1} \mathbf{S}(\theta_0) $$
Under $H_0$, $T_S \xrightarrow{d} \chi^2_p$.
We reject for large values of $T_S$.

\subsection{Lagrange Multiplier Interpretation}
When testing composite hypotheses like $g(\theta) = 0$, we maximize the likelihood subject to constraints.
The Score Test is mathematically equivalent to the \textbf{Lagrange Multiplier Test}.
If we maximize $\ell(\theta) - \lambda^T g(\theta)$, the value of the Lagrange multiplier $\lambda$ measures how much the constraint "hurts" the likelihood.
\begin{itemize}
    \item If $\lambda \approx 0$, the constraint is consistent with data ($H_0$ likely true).
    \item If $\lambda$ is large, the constraint forces the solution far from the unconstrained peak ($H_0$ likely false).
\end{itemize}
Rao's statistic effectively measures the magnitude of this constraint force.

% -------------------------------------------------------------------
% PAGE 303 - 357: Confidence Intervals & Multiple Testing
% -------------------------------------------------------------------
\section{Interval Estimation}

In point estimation, we acknowledge variability by computing standard errors. In a similar spirit, instead of a single value, we may construct an interval that contains the true parameter with high probability.

\subsection{Definitions}
\begin{definition}[\textbf{Confidence Interval}]
Let $\Xvec$ be a random sample. A confidence interval is a random interval of the form $[L(\Xvec), U(\Xvec)]$, where $L$ and $U$ are statistics satisfying $L(\Xvec) \le U(\Xvec)$ almost surely.
\end{definition}

\begin{definition}[\textbf{Confidence Set}]
More generally, a random set $S(\Xvec) \subset \Theta$ is called a confidence set.
\end{definition}

\begin{example}[\textbf{Normal Confidence Set}]
Let $X_1, \dots, X_n \iid N(\mu, \sigma^2)$.
We might define a set for the vector $(\mu, \sigma^2)$:
$$ S(\Xvec) = \{ (\mu, \sigma^2) : \bar{X} - 2\sigma \le \mu \le \bar{X} + 2\sigma, \sigma^2 > 0 \} $$
\end{example}

\subsection{Coverage Probability}
We want the interval to contain the true parameter $\theta$.
The Coverage Probability for a specific $\theta$ is:
$$ P_\theta(\theta \in S(\Xvec)) = P_\theta(L(\Xvec) \le \theta \le U(\Xvec)) $$
The Confidence Coefficient is the infimum of the coverage probability over the parameter space:
$$ \text{Confidence Coefficient} = \inf_{\theta \in \Theta} P_\theta(\theta \in S(\Xvec)) $$

\textbf{Note:} The probability statement is about the random interval $S(\Xvec)$, not the fixed parameter $\theta$.

\section{Construction of Confidence Intervals}

There are two primary methods:
1. Using Test Statistics (Inverting Tests).
2. Using Pivotal Quantities.

\subsection{Using Test Statistics}
Consider $X_i \iid N(\mu, 1)$. We want a CI for $\mu$.
Recall the test for $H_0: \mu = \mu_0$ uses $Z = \sqrt{n}(\bar{X} - \mu_0)$.
We know $Z \sim N(0, 1)$.
$$ P(-z_{\alpha/2} \le \sqrt{n}(\bar{X} - \mu) \le z_{\alpha/2}) = 1 - \alpha $$
Rearranging the inequality for $\mu$:
$$ \bar{X} - \frac{z_{\alpha/2}}{\sqrt{n}} \le \mu \le \bar{X} + \frac{z_{\alpha/2}}{\sqrt{n}} $$
This random interval has coverage probability $1-\alpha$ for every $\mu$. Thus, it is a $100(1-\alpha)\%$ Confidence Interval.

\subsection{Using Pivots}
A Pivotal Quantity (or Pivot) is a function $Q(\Xvec, \theta)$ whose distribution is free of $\theta$.
Example: $Q = \frac{\bar{X} - \mu}{S/\sqrt{n}} \sim t_{n-1}$.
Since the distribution is known (t-distribution) and does not depend on $\mu$ or $\sigma$, we can fix probabilities:
$$ P(-t_{\alpha/2} \le Q \le t_{\alpha/2}) = 1 - \alpha $$
Inverting this for $\mu$ yields the standard t-interval.

\section{Duality: Tests and Intervals}

There is a one-to-one correspondence between acceptance regions of hypothesis tests and confidence sets.

\begin{theorem}[\textbf{Test Inversion}]
Let $\phi_{\theta_0}$ be a level $\alpha$ test for $H_0: \theta = \theta_0$.
Let $\mathcal{A}(\theta_0)$ be the acceptance region for this test.
Define the set $S(\xvec)$ as all parameter values $\theta$ for which the test accepts given data $\xvec$:
$$ S(\xvec) = \{ \theta \in \Theta : \xvec \in \mathcal{A}(\theta) \} $$
Then $S(\Xvec)$ is a confidence set with confidence coefficient $1 - \alpha$.
\end{theorem}

\subsection{One-Sided Intervals}
We can obtain one-sided intervals by inverting one-sided tests.
For $N(\mu, 1)$, testing $H_0: \mu = \mu_0$ vs $H_1: \mu > \mu_0$.
The acceptance region is $\sqrt{n}(\bar{X} - \mu_0) \le z_\alpha$ (reject for large values).
Inverting for $\mu_0$:
$$ \mu_0 \ge \bar{X} - \frac{z_\alpha}{\sqrt{n}} $$
This gives the one-sided lower confidence bound:
$$ \left[ \bar{X} - \frac{z_\alpha}{\sqrt{n}}, \infty \right) $$

\textbf{Proof of Duality:}
$$ P_\theta(\theta \in S(\Xvec)) = P_\theta(\Xvec \in \mathcal{A}(\theta)) $$
Since $\mathcal{A}(\theta)$ is the acceptance region of a size $\alpha$ test for null value $\theta$, this probability is exactly $1 - \alpha$.

\begin{example}[\textbf{Normal Mean Inversion}]
$X_i \iid N(\mu, 1)$. Test $\phi(\xvec)$ rejects if $|\sqrt{n}(\bar{X} - \mu_0)| > z_{\alpha/2}$.
Acceptance region $\mathcal{A}(\mu_0)$:
$$ \mathcal{A}(\mu_0) = \{ \xvec : -z_{\alpha/2} \le \sqrt{n}(\bar{X} - \mu_0) \le z_{\alpha/2} \} $$
Inverting for $\mu_0$:
$$ S(\xvec) = \{ \mu : \bar{X} - \frac{z_{\alpha/2}}{\sqrt{n}} \le \mu \le \bar{X} + \frac{z_{\alpha/2}}{\sqrt{n}} \} $$
\end{example}

\subsection{Optimality of Intervals}
\begin{itemize}
    \item \textbf{Length:} Shorter intervals are better (more precise).
    \item \textbf{Relation to UMP:} If we start with a UMP (or UMP Unbiased) test, the corresponding inverted confidence set usually has the shortest expected length (or smallest volume).
    \item Example: The UMPU test for the normal mean leads to the standard symmetric interval, which is the shortest among all unbiased intervals.
\end{itemize}

\subsection{Approximate Intervals (Wald)}
Often exact pivots are unavailable. We use asymptotic pivots from Wald's test.
Recall $T_W = (\hat{\theta} - \theta)^2 I(\hat{\theta})$.
$$ (\hat{\theta} - \theta) \sqrt{I(\hat{\theta})} \xrightarrow{d} N(0, 1) $$
An approximate $100(1-\alpha)\%$ CI is:
$$ \hat{\theta} \pm z_{\alpha/2} \frac{1}{\sqrt{I(\hat{\theta})}} $$
This is the interval obtained by inverting the Wald test.
We cannot guarantee exactly $1-\alpha$ coverage for small $n$, only asymptotically.

\section{Multiparameter Confidence Sets}

Let $\theta \in \R^d$. We want a region $S(\Xvec) \subset \R^d$.
We can invert multiparameter tests.

\subsection{1. Inverting LRT (Likelihood Regions)}
The acceptance region for LRT is $\{ \xvec : -2 \ln \lambda(\xvec) \le \chi^2_{d, \alpha} \}$.
The confidence set is:
$$ S(\xvec) = \{ \theta : 2[\ell(\hat{\theta}_{MLE}) - \ell(\theta)] \le \chi^2_{d, \alpha} \} $$
This defines a contour of the log-likelihood function. It is often not an ellipsoid but creates valid regions that respect the boundaries of the parameter space.

\subsection{2. Inverting Wald (Ellipsoids)}
Wald statistic: $T_W = (\hat{\theta} - \theta)^T \widehat{\mathbf{I}} (\hat{\theta} - \theta)$.
Acceptance region: $T_W \le \chi^2_{d, \alpha}$.
Confidence set:
$$ S(\xvec) = \{ \theta : (\theta - \hat{\theta})^T \widehat{\mathbf{I}} (\theta - \hat{\theta}) \le \chi^2_{d, \alpha} \} $$
This equation defines an ellipsoid centered at the MLE $\hat{\theta}$. The shape is determined by the Fisher Information matrix.

\subsection{3. Inverting Score Test}
Rao's Score statistic: $T_S = S(\theta)^T [\mathbf{I}(\theta)]^{-1} S(\theta)$.
Confidence set:
$$ S(\xvec) = \{ \theta : S(\theta)^T [\mathbf{I}(\theta)]^{-1} S(\theta) \le \chi^2_{d, \alpha} \} $$
This method is useful because it does not require computing the MLE for every potential $\theta$ inside the set, though solving the inequality for $\theta$ can be numerically complex.

\section{Multiple Hypothesis Testing}

Consider the problem of testing multiple hypotheses simultaneously.
\begin{example}
We want to study the effect of a drug on several vitals (blood sugar, blood pressure, heart rate, etc.).
Data: $X_{ij}$ is the difference for the $i$-th person on the $j$-th vital.
$$ X_{1j}, \dots, X_{nj} \iid N(\mu_j, \sigma_j^2) $$
If $\mu_j = 0$, the drug has no effect on the $j$-th vital.
We want to test $H_{0j}: \mu_j = 0$ vs $H_{1j}: \mu_j \neq 0$ for $j=1, \dots, d$.
\end{example}

We face two questions:
\begin{enumerate}
    \item Does the drug affect \textit{any} vital? (Global Null $H_0 = \bigcap H_{0j}$)
    \item Which specific vitals are affected?
\end{enumerate}

\subsection{The Problem of Multiplicity}
Suppose we test each component hypothesis $H_{0j}$ independently at level $\alpha$.
If the vitals are independent, the probability of at least one false rejection (Type I Error for the global null) is:
$$ P(\text{Reject at least one } H_{0j} | \text{All } H_{0j} \text{ are true}) = 1 - (1-\alpha)^d $$
For $d=10$ and $\alpha=0.05$, this probability is $1 - (0.95)^{10} \approx 0.40$.
For $d=100$, it is $\approx 0.99$.
Thus, testing individually without correction leads to a very high chance of finding spurious effects.

\subsection{Bonferroni Correction}
To control the global error rate at $\alpha$, we can perform each component test at level $\alpha/d$.
$$ P(\bigcup_{j=1}^d \{ \text{Reject } H_{0j} \}) \le \sum_{j=1}^d P(\text{Reject } H_{0j}) = \sum \frac{\alpha}{d} = \alpha $$
(Using Boole's Inequality/Union Bound).
This works regardless of dependence between vitals. However, it is very conservative (low power).

\section{P-Values in Multiple Testing}

Methods are usually defined in terms of $p$-values.
Let $P_j$ be the $p$-value for hypothesis $H_{0j}$.
Under the null $H_{0j}$, $P_j \sim \text{Uniform}(0, 1)$.
Sorted p-values: $P_{(1)} \le P_{(2)} \le \dots \le P_{(d)}$.
Corresponding Hypotheses: $H_{(1)}, H_{(2)}, \dots, H_{(d)}$.

\textbf{Bonferroni Rule:} Reject $H_{0j}$ if $P_j \le \alpha/d$.

\section{Error Metrics}

When testing $d$ hypotheses, we can classify our decisions:

\begin{table}[h]
\centering
\begin{tabular}{|c|c|c|c|}
\hline
& \textbf{Accept Null} & \textbf{Reject Null} & \textbf{Total} \\
\hline
\textbf{Null True} & $U$ & $V$ (False Positives) & $m_0$ \\
\hline
\textbf{Null False} & $T$ (False Negatives) & $S$ (True Positives) & $d - m_0$ \\
\hline
\textbf{Total} & $d - R$ & $R$ & $d$ \\
\hline
\end{tabular}
\end{table}

Here $m_0$ is the number of true null hypotheses (unknown). $R$ is the number of rejections (observable). $V$ is the unobservable random variable representing the number of False Rejections.

\subsection{Error Rates}
\begin{itemize}
    \item \textbf{FWER (Family-Wise Error Rate):} The probability of making at least one false rejection.
    $$ \text{FWER} = P(V \ge 1) $$
    \item \textbf{PCE (Per Comparison Error):} Expected error per test. $E[V]/d$.
    \item \textbf{PFE (Per Family Error):} Expected number of false rejections. $E[V]$.
\end{itemize}
Relationship: $\text{PCE} \le \text{FWER} \le \text{PFE}$.
Bonferroni controls FWER.

\subsection{False Discovery Rate (FDR)}
In modern applications (e.g., genomics with thousands of tests), FWER is too strict. We are willing to accept some false positives to find more true signals.
We define the False Discovery Proportion (FDP):
$$ Q = \begin{cases} V/R & \text{if } R > 0 \\ 0 & \text{if } R = 0 \end{cases} $$
The False Discovery Rate (FDR) is the expected FDP:
$$ \text{FDR} = \E[Q] $$

\section{Stepwise Procedures}

\subsection{1. Holm's Method (Step-Down)}
An improvement over Bonferroni that also controls FWER.
\begin{enumerate}
    \item Order p-values: $P_{(1)} \le P_{(2)} \le \dots \le P_{(d)}$.
    \item Step 1: Check if $P_{(1)} \le \alpha/d$. If not, accept all. Else reject $H_{(1)}$ and continue.
    \item Step $k$: Check if $P_{(k)} \le \alpha/(d-k+1)$. If not, accept $H_{(k)}, \dots, H_{(d)}$. Else reject $H_{(k)}$ and continue.
\end{enumerate}
This is uniformly more powerful than Bonferroni.

\subsection{2. Benjamini-Hochberg (BH) Procedure}
This procedure controls the FDR at level $\alpha$ (assuming independence or positive dependence).
\begin{enumerate}
    \item Order p-values: $P_{(1)} \le \dots \le P_{(d)}$.
    \item Find the largest index $k$ such that:
    $$ P_{(k)} \le \frac{k}{d} \alpha $$
    \item Reject all hypotheses $H_{(1)}, \dots, H_{(k)}$.
\end{enumerate}
If no such $k$ exists, reject nothing.

\section{U-Statistics}

Statistics formed by averaging a kernel over the sample.
Let $X_1, \dots, X_n$ be a sample.
Let $h(x_1, \dots, x_k)$ be a \textbf{symmetric kernel} function of degree $k$.

The U-statistic corresponding to $h$ is:
$$ U(\Xvec) = \frac{1}{\binom{n}{k}} \sum_{\beta} h(X_{\beta_1}, \dots, X_{\beta_k}) $$
where the sum is over all distinct combinations (subsets) of size $k$ from $\{1, \dots, n\}$.

\begin{example}[\textbf{Sample Mean}]
$k=1$. $h(x) = x$.
$$ U = \frac{1}{n} \sum X_i = \bar{X} $$
\end{example}

\begin{example}[\textbf{Sample Variance}]
$k=2$. $h(x_1, x_2) = \frac{1}{2}(x_1 - x_2)^2$.
This leads to the unbiased sample variance $S^2 = \frac{1}{n-1} \sum (X_i - \bar{X})^2$.
\end{example}

\textbf{Properties:}
\begin{itemize}
    \item Unbiasedness: $\E[U] = \theta$ where $\theta = \E[h(X_1, \dots, X_k)]$.
    \item Variance: Given by $\Var(U) = \binom{n}{k}^{-1} \sum_{c=1}^k \binom{k}{c} \binom{n-k}{k-c} \zeta_c$ where $\zeta_c = \Var(\E[h | X_1, \dots, X_c])$.
\end{itemize}
Other examples include Kendall's Tau and Spearman's Rho.

\end{document}